% Chapter 1: Why Soft Cohesion? Limits of Traditional Approaches
% Part I: Foundations
% Expanded edition: ~27pp, 8 figures

\chapter{왜 Soft Cohesion인가?\\전통적 접근의 한계}
\label{ch:motivation}

\begin{quote}
\textit{``물건은 어디서 시작되고, 어디서 끝나는가?''}
\end{quote}

\section{도입: 인식의 기본 문제}
\label{sec:intro-problem}

우리는 세상을 ``물건들''의 집합으로 경험한다. 책상 위의 컵, 창밖의 나무, 화면 위의 글자---이들은 각각 독립적이고 경계가 뚜렷한 개체로 인식된다. 이러한 인식은 너무나 자연스럽기 때문에 우리는 좀처럼 그것이 얼마나 놀라운 성취인지를 깨닫지 못한다.

그러나 조금만 주의를 기울이면, ``물건''의 경계가 그리 자명하지 않은 경우가 수없이 많음을 알게 된다:

\begin{itemize}
    \item \textbf{안개 속의 산}: 산의 윤곽은 안개 속에서 점차 사라진다. ``산''은 어디서 끝나고 ``안개''는 어디서 시작되는가?
    \item \textbf{군중 속의 무리}: 축제 현장에서 사람들이 자연스럽게 모여든 무리는 명확한 경계가 없다. 가장자리에 있는 사람은 무리의 일부인가, 아닌가?
    \item \textbf{새벽의 빛}: 밤에서 낮으로의 전이는 ``어둠''과 ``밝음''의 경계가 존재하지 않는 연속적 과정이다.
    \item \textbf{수채화의 번짐}: 물감이 종이 위에서 퍼져나갈 때, 한 색의 영역과 다른 색의 영역 사이에는 점진적 전이만이 있다.
    \item \textbf{음악의 화성}: 오케스트라에서 하나의 ``소리''는 여러 악기의 음이 응집하여 형성된다. 각 악기의 기여는 연속적이며, 그 경계는 모호하다.
\end{itemize}

이 예들이 보여주는 것은, 물건의 경계가 선명한 것은 예외이지 규칙이 아니라는 점이다. 대부분의 자연적 형성(formation)은 \emph{정도(degree)}를 가진다---완전한 소속과 완전한 비소속 사이에는 연속적 전이가 존재한다.

그렇다면 ``물건''이란 실제로 무엇인가? 보다 근본적으로, 물건이 하나의 통일된 개체로 \emph{형성}되는 과정은 어떻게 일어나는가? 이 질문은 인지과학, 컴퓨터 비전, 수리 형태학의 교차점에 있는 근본적 문제이다.

전통적 접근법은 이 질문을 한 방향으로만 접근해왔다: 먼저 ``물건''이 있다고 전제한 뒤, 그것을 감지하거나 분류하거나 추적하는 알고리즘을 설계하는 것이다. 그러나 이러한 접근은 핵심 문제를 회피한다---물건이 아직 물건으로 존재하기 \emph{이전}의 상태, 즉 개별화(individuation)가 일어나기 전의 세계를 기술하지 못한다.

\textbf{Soft Cognitive Cohesion}(부드러운 인지 응집, SCC) 이론은 이 문제를 정면으로 다룬다. SCC는 물건을 출발점으로 삼지 않는다. 대신 \emph{응집(cohesion)}---관계적 자기-지지(relational self-support)의 연속적 장(field)---을 기본 구조로 삼고, 물건은 이 장이 특정 조건을 충분히 만족할 때 사후적으로 읽어낼 수 있는 안정화된 패턴으로 취급한다.

이 장에서는 왜 이러한 전환이 필요한지, 전통적 접근이 어디서 한계에 부딪히는지, 그리고 SCC가 어떤 대안을 제시하는지를 살펴본다.


\section{전통적 객체 중심 접근의 문제}
\label{sec:traditional-limits}

전통적 접근법은 세 가지 핵심적인 가정 위에 서 있다:
\begin{enumerate}
    \item 세계는 이산적(discrete) 개체들의 집합으로 주어진다.
    \item 각 개체는 명확한 경계를 가진다.
    \item 시간에 걸쳐 각 개체는 고유한 정체성(identity)을 유지한다.
\end{enumerate}

이 가정들 각각에 심각한 문제가 있다. 표~\ref{tab:framework-comparison}은 주요 프레임워크들이 이 가정에 어떻게 의존하는지를 정리한다.

\begin{table}[t]
\centering
\caption{기존 프레임워크의 암묵적 가정 비교}
\label{tab:framework-comparison}
\begin{tabular}{p{2.8cm}ccccp{3cm}}
\hline
\textbf{프레임워크} & \textbf{이산 개체} & \textbf{명확 경계} & \textbf{고유 ID} & \textbf{전-객체} & \textbf{핵심 한계} \\
\hline
영상 분할 & 전제 & 전제 & --- & 불가 & 경계가 이진적 \\
객체 추적 & 전제 & 전제 & 전제 & 불가 & 분열/합병 처리 불가 \\
클러스터링 & 암묵 전제 & 파생 & --- & 불가 & 거리 기반, 자기-참조 없음 \\
퍼지 집합 & --- & 완화 & --- & 부분적 & 형성 역학 없음 \\
위상 장 이론 & --- & 완화 & --- & 부분적 & 단일 에너지, 자기-참조 없음 \\
\hline
\textbf{SCC} & 파생 & 파생 & 파생 & \textbf{가능} & (이 책의 주제) \\
\hline
\end{tabular}
\end{table}


\subsection{분할의 문제: 경계는 어디에 있는가?}
\label{subsec:segmentation-problem}

영상 분할(image segmentation)은 컴퓨터 비전의 가장 기본적 과제 중 하나이다. 입력 영상을 ``물건''에 해당하는 영역들로 나누는 것이 목표인데, 이는 경계가 항상 명확하다는 전제 위에 서 있다.

그러나 현실의 경계는 대부분 \emph{점진적}이다. 안개 속의 산, 수면에 비친 반영, 붓으로 번진 수채화---이들의 경계는 이진적(binary)이지 않다. 소속과 비소속 사이에는 연속적인 전이 대역(transition band)이 존재한다. 전통적 분할은 이 전이 대역을 억지로 칼로 자르듯 이분법적으로 나눈다:
\begin{equation}
    A_t = \{x \in X : f(x) \geq \theta\}
    \label{eq:crisp-threshold}
\end{equation}
여기서 $\theta$는 분할 문턱값이다. 이러한 이분법은 전이 대역의 풍부한 구조적 정보를 파괴한다.

더 근본적으로, 분할은 분할 \emph{이전에} 이미 ``무엇이 하나의 물건인가''가 결정되어 있어야 한다. 그러나 이 결정 자체가 분할이 해결하려는 문제이므로, 논리적 순환이 발생한다.

% [Figure placeholder: crisp segmentation vs. graded field]
\begin{figure}[t]
\centering
% \includegraphics[width=0.9\textwidth]{figures/ch01_crisp_vs_graded.pdf}
\fbox{\parbox{0.85\textwidth}{\centering\vspace{4cm}
\textbf{[그림 1.1]} 좌: 전통적 이진 분할---명확한 경계선, 전이 대역 소실.\\
우: 연속적 응집 장(cohesion field)---내부(1에 가까움), 전이 대역(중간값), 외부(0에 가까움)의 풍부한 구조 보존.
\vspace{0.5cm}}}
\caption{객체 중심 접근(좌)과 응집 장 접근(우)의 비교. 전통적 분할은 임계값 $\theta$에서 날카로운 경계를 강제하지만, 응집 장은 내부에서 외부로의 점진적 전이를 자연스럽게 표현한다.}
\label{fig:crisp-vs-graded}
\end{figure}


\subsection{추적의 문제: 정체성은 무엇인가?}
\label{subsec:tracking-problem}

객체 추적(object tracking)은 시간에 걸쳐 ``같은 물건''을 식별하는 과제이다. 전형적 추적 시스템은 각 개체에 고유 식별자(tracking ID)를 부여하고, 이를 프레임 간에 대응시킨다.

그러나 ``같은 물건''이라는 개념 자체가 문제적이다. 물건이 분열하면? 합쳐지면? 점진적으로 변형되어 원래와 전혀 다른 모습이 되면? 추적 ID 기반의 접근은 이러한 경우에 본질적 한계를 드러낸다. 정체성을 미리 전제하기 때문에, 정체성이 \emph{형성}되는 과정을 기술할 수 없다.

SCC 이론에서 시간적 정체성은 전혀 다르게 정의된다. 정체성은 추적 ID가 아니라 \textbf{시간적 전달 핵(temporal transport kernel)} $\mathbf{M}_{t \to s}$에 의한 구조적 계승이다. 형성의 응집적 핵심이 시간 전달을 통해 계승될 때, 그 형성은 \emph{지속(persist)}한다. 이는 철학적 발생적 정체성(genidentity) 개념과 깊이 연결된다.


\subsection{경계의 문제: 내부와 외부}
\label{subsec:boundary-problem}

전통적 접근에서 물건의 경계는 날카로운 선(sharp line)이다---소속 집합의 위상적 경계(topological boundary). 그러나 전-객체적(pre-objective) 단계에서 경계는 날카로운 선이 아니라 \textbf{전이 대역(transition band)}이다:
\begin{equation}
    \mathrm{Bd}_t(u_t) = \{x \in X_t \mid \theta_1 < u_t(x) < \theta_2\}
    \label{eq:boundary-band}
\end{equation}
이 대역 내에서 응집은 완전히 확립된 것도 아니고 완전히 부재한 것도 아니다. 최대한의 구조적 전이가 일어나는 영역---형성의 내부적 관계 지지가 외부로 소실되는 곳---이 바로 경계 대역이다.

이러한 점진적 경계 개념은 이진 분할로는 원천적으로 포착할 수 없다.


\section{전-객체적 이론은 어떤 모습이어야 하는가?}
\label{sec:desiderata}

위에서 살펴본 한계들은 공통의 근원을 가진다: 기존 접근법은 개별화가 이미 완료된 상태를 전제한다. 그렇다면, 개별화 \emph{이전}의 상태를 기술할 수 있는 이론은 어떤 성질을 가져야 하는가? 다섯 가지 필요조건(desiderata)을 제시한다.

\textbf{D1. 연속적 장(Graded Fields).}
이론의 기본 존재자는 이진적 집합이 아니라 $[0,1]$-값의 연속적 장이어야 한다. 소속의 정도를 허용해야 ``부분적으로 형성된'' 상태를 기술할 수 있다.

% [Figure placeholder: D1 - graded vs. binary]
\begin{figure}[t]
\centering
\fbox{\parbox{0.85\textwidth}{\centering\vspace{3cm}
\textbf{[그림 1.2]} 다섯 가지 필요조건의 시각적 예시.\\
D1: 이진 집합 vs.\ 연속 장.\ D2: 자기-참조적 순환 구조.\\
D3: 점진적 내외부 전이.\ D4: 전달 핵에 의한 구조적 계승.\\
D5: 에너지 최소화 경관.
\vspace{0.5cm}}}
\caption{전-객체적 이론의 다섯 가지 필요조건. 각 조건은 기존 프레임워크에서 충족되지 않는 고유한 요구를 표현한다.}
\label{fig:five-desiderata}
\end{figure}

\textbf{D2. 자기-참조적 구조(Self-Referential Structure).}
장의 ``좋은 형태''를 판단하는 기준이 장 자체로부터 파생되어야 한다. 외부에서 주어진 기준(클래스 라벨, 거리 임계값)에 의존해서는 안 된다. 형성은 자기 자신의 관계 구조에 비추어 스스로를 평가할 수 있어야 한다.

\textbf{D3. 내부/외부의 창발(Emergence of Inside/Outside).}
내부와 외부의 구분은 이론의 출발점이 아니라 \emph{결과}여야 한다. 연속적 장에서 내부 영역(높은 응집)과 외부 영역(낮은 응집) 사이의 비대칭이 형성 과정을 통해 자연스럽게 발생해야 한다.

\textbf{D4. 추적 ID 없는 시간적 계승(Temporal Inheritance Without Tracking IDs).}
시간적 정체성은 고유 식별자의 대응이 아니라 구조적 조직의 계승으로 정의되어야 한다. 형성이 이동하고, 변형되고, 부분적으로 소실되는 상황을 자연스럽게 다룰 수 있어야 한다.

\textbf{D5. 에너지 기반 특성화(Energy-Based Characterization).}
형성은 규칙(rule)이 아니라 에너지 원리에 의해 특성화되어야 한다. 에너지 최소화는 형성의 존재, 유일성, 안정성, 수렴에 대한 수학적 도구를 제공한다. 이는 ``형성은 왜 그 형태인가?''라는 질문에 답할 수 있게 해준다.

이 다섯 조건을 모두 충족하는 기존 프레임워크는 없다 (표~\ref{tab:framework-comparison} 참조). SCC 이론은 이 다섯 조건의 동시 충족을 목표로 설계된다.


\section{대안: 응집도(Cohesion) 우선}
\label{sec:cohesion-first}

전통적 접근의 공통된 한계는 ``너무 늦게 시작한다''는 것이다. 이미 개별화가 완료된 상태를 전제하고, 그 결과물을 다루려 한다. SCC 이론은 이 순서를 뒤집는다.

\begin{center}
\fbox{\parbox{0.8\textwidth}{\centering
\textbf{존재론적 우선성 원리 (Ontological Priority Principle)}\\[6pt]
부드러운 응집 장(soft cohesion field) $u_t : X_t \to [0,1]$은\\
이론의 \emph{원초적 존재자}(primitive entity)이다.\\
이산적 객체는 항상 파생적이다---이 우선성을 결코 역전시키지 않는다.
}}
\end{center}

이 원리가 의미하는 바를 구체적으로 살펴보자.

\subsection{응집 장의 의미}

응집 장 $u_t(x) \in [0,1]$에서 각 값은 해당 지점 $x$가 응집적 형성에 참여하는 \emph{강도}를 나타낸다. 이것은 확률이 아니다. 클래스 소속 점수도 아니다. 지점 $x$가 형성의 내부적 관계 조직에 의해 얼마나 강하게 유지되고(sustained), 동시에 그 조직에 얼마나 기여하는지(contributes)를 표현하는 \textbf{응집적 참여 강도}이다.

\begin{itemize}
    \item $u_t(x) \approx 1$: 깊은 내부성(deep interiority)---강하고 안정적으로 형성에 통합됨
    \item $u_t(x) \approx 0$: 실질적 외부성(effective exteriority)---형성에 참여하지 않음
    \item $0 < u_t(x) < 1$: 전이적 또는 경계적 참여---소속과 비소속 사이의 구조적 전이
\end{itemize}

\subsection{왜 부드러운 형태가 원초적인가}

부드러운(soft) 체계가 이산적(crisp) 체계보다 원초적인 이유는 네 가지이다.

\textbf{(1) 전-객체적 응집은 본질적으로 연속적이다.} 형성이 인식 가능한 객체로 안정화되기 이전에, 응집은 정도를 허용한다. 어떤 영역은 강하게, 어떤 영역은 약하게, 어떤 영역은 모호하게 참여한다. 이진 집합 $A_t \subseteq X_t$는 이 연속적 전-객체 단계를 즉각 왜곡한다.

\textbf{(2) 내부와 외부는 처음부터 날카롭게 분리되지 않는다.} 형성의 내부와 외부 사이에는 전형적으로 전이 대역이 존재한다. 이진 집합에는 이러한 전이 대역이 없다.

\textbf{(3) 닫힘(closure)은 완성된 것이 아니라 형성 중인 것이다.} 이산적 경우에 닫힘 연산은 보통 멱등(idempotent)이다: 두 번 적용해도 한 번과 같다. 그러나 전-객체 수준에서 닫힘은 안정화 \emph{경향}이지, 이미 완료된 연산이 아니다. 부드러운 체계는 이를 자연스럽게 수용한다.

\textbf{(4) 이산적 체계는 부드러운 체계에서 복원 가능하다.} 응집 장 $u_t$가 주어지면, 임계값을 적용하여 이산 부분집합을 얻을 수 있다: $A_t = \{x \mid u_t(x) \geq \theta\}$. 이산 경계, 이산 내부, 이산 정체성 판정도 모두 복원 가능하다. 그러나 이들은 더 풍부한 부드러운 구조의 \emph{사영(projection)}이지, 그 역은 아니다.


\section{SCC의 핵심 아이디어}
\label{sec:scc-core-ideas}

SCC 이론은 다음과 같은 핵심 구조 위에 세워진다.


\subsection{응집 장의 세 가지 예}

SCC의 핵심 개념을 직관적으로 이해하기 위해, $10 \times 10$ 격자 위의 세 가지 응집 장을 비교해보자.

% [Figure placeholder: three field examples]
\begin{figure}[t]
\centering
\fbox{\parbox{0.85\textwidth}{\centering\vspace{4cm}
\textbf{[그림 1.3]} $10 \times 10$ 격자 위의 세 가지 응집 장.\\
좌: Gaussian bump---잘 형성된 단일 형성 ($\mathsf{Bind} \approx 0.97$, $\mathsf{Sep} \approx 0.85$).\\
중앙: 균일 장 $u \equiv 0.3$---높은 Bind이지만 Sep$\approx 0.5$ (형성 없음).\\
우: 무작위 장---모든 진단 값이 낮음 (구조 없음).
\vspace{0.5cm}}}
\caption{응집 장의 세 가지 전형적 유형. 색상은 응집 강도를 나타낸다 (빨강=높음, 파랑=낮음). Gaussian bump만이 진정한 형성이며, 균일 장과 무작위 장은 각각 다른 이유로 형성이 아니다.}
\label{fig:three-fields}
\end{figure}

\textbf{(1) Gaussian bump}: 격자 중앙에 집중된 원형 형성. 내부(높은 값)와 외부(낮은 값) 사이에 자연스러운 전이 대역이 있다. 네 가지 진단 요소 모두 높은 값을 가진다---이것이 ``좋은 형성''이다.

\textbf{(2) 균일 장} $u \equiv c$: 모든 지점에서 동일한 값. 관계적으로는 완벽히 자기-지지적이지만(Bind $\approx$ 1), 내외부 구분이 전혀 없다(Sep $\approx$ 0.5). 형태학적 구조도 없다(Inside = 0). 이것은 형성이 \emph{아니다}---Bind만으로는 충분하지 않음을 보여준다.

\textbf{(3) 무작위 장}: 독립 균등 분포에서 생성된 값. 관계적 지지가 없고(낮은 Bind), 구조가 없으며(낮은 Inside), 외부와 구분되지 않는다(중간 Sep). 모든 진단 요소가 낮다---이것은 어떤 의미에서도 형성이 아니다.

이 비교가 보여주는 핵심 교훈: 형성의 질은 \emph{네 가지 독립적 차원}에서 동시에 평가되어야 한다. 하나의 숫자로는 부족하다.


\subsection{형식 우주}

이론의 형식 우주는 구조화된 튜플이다:
\begin{equation}
    \mathfrak{C}^{\mathrm{soft}} = \Big( T,\ \{X_t\}_{t \in T},\ \{u_t\}_{t \in T},\ \{\mathrm{Cl}_t\}_{t \in T},\ \{\mathbf{N}_t, \mathbf{D}_t\}_{t \in T},\ \{\mathbf{M}_{t \to s}\}_{t,s \in T} \Big)
    \label{eq:formal-universe}
\end{equation}

각 구성 요소의 직관적 의미는 다음과 같다:
\begin{itemize}
    \item $T$: 시간 지표 집합---시간 순서가 잘 정의된 전순서 집합
    \item $X_t$: 시각 $t$에서의 관계적 지지 공간---응집이 정의되는 장소들의 집합
    \item $u_t : X_t \to [0,1]$: 부드러운 응집 장---이론의 원초적 존재자
    \item $\mathrm{Cl}_t$: 부드러운 닫힘 연산자---관계적 자기-지지의 완성
    \item $\mathbf{N}_t$: 부드러운 인접 핵---국소적 관계 구조
    \item $\mathbf{D}_t$: 부드러운 구별 연산자---내부와 외부의 비대칭
    \item $\mathbf{M}_{t \to s}$: 시간적 전달 핵---응집 조직의 시간적 계승
\end{itemize}


\subsection{연산자 쌍: 자기-완성과 자기-대비}

SCC의 내부 역학은 두 연산자의 상호작용으로 구동된다:

\textbf{자기-완성(Self-completion)}: 닫힘 연산자 $\mathrm{Cl}_t$는 응집 장을 그 관계적으로 완성된 형태로 매핑한다. 직관적으로, 각 지점의 응집도를 그 이웃들로부터 받는 관계적 지지에 맞게 갱신하는 것이다.

\textbf{자기-대비(Self-contrast)}: 구별 연산자 $\mathbf{D}_t$는 각 지점이 외부 장 $1-u_t$에 대해 가지는 구조적 비대칭을 측정한다. 직관적으로, 내부와 외부가 얼마나 구분되는지를 평가하는 것이다.

이 두 연산자의 상호작용이 형성의 자기-조직화를 구동한다: 닫힘은 내부 응집을 강화하고, 구별은 내외부 분리를 강화한다.


\subsection{에너지 원리}

SCC는 변분적(variational) 이론이다. 형성은 네 항의 에너지 범함수를 최소화하는 응집 장으로 특징지어진다:
\begin{equation}
    \mathcal{E}(\mathbf{u}) = \lambda_{\mathrm{cl}}\, \mathcal{E}_{\mathrm{cl}} + \lambda_{\mathrm{sep}}\, \mathcal{E}_{\mathrm{sep}} + \lambda_{\mathrm{bd}}\, \mathcal{E}_{\mathrm{bd}} + \lambda_{\mathrm{tr}}\, \mathcal{E}_{\mathrm{tr}}
    \label{eq:energy-preview}
\end{equation}
네 항은 각각 독립적인 구조적 요구를 반영한다:
\begin{enumerate}
    \item $\mathcal{E}_{\mathrm{cl}}$ (닫힘): 내부적 관계 자기-지지
    \item $\mathcal{E}_{\mathrm{sep}}$ (분리): 외부와의 구조적 대비
    \item $\mathcal{E}_{\mathrm{bd}}$ (경계/형태학): 공간적 분절의 일관성
    \item $\mathcal{E}_{\mathrm{tr}}$ (전달): 시간적 응집 조직의 계승
\end{enumerate}

이 네 항의 독립성은 이론의 근본적 약속이다: 어느 것도 다른 것으로 환원되지 않으며, 각각은 형성이 실패할 수 있는 독립적 차원을 포착한다. 자세한 수학적 정의는 제~\ref{ch:energy}장에서 다룬다.


% [Figure placeholder: theory pipeline]
\begin{figure}[p]
\centering
\fbox{\parbox{0.9\textwidth}{\centering\vspace{10cm}
\textbf{[그림 1.4]} SCC 이론의 전체 파이프라인 (전면 페이지).\\[8pt]
\textbf{존재론 (형식 우주):} $T, X_t, u_t, \mathrm{Cl}_t, \mathbf{N}_t/\mathbf{D}_t, \mathbf{M}_{t \to s}$\\[4pt]
$\downarrow$\\[4pt]
\textbf{공리론 (5개 그룹):} A (닫힘), B (인접), C (공-소속), D (구별), E (전달)\\[4pt]
$\downarrow$\\[4pt]
\textbf{에너지 원리:} $\mathcal{E} = \lambda_{\mathrm{cl}} \mathcal{E}_{\mathrm{cl}} + \lambda_{\mathrm{sep}} \mathcal{E}_{\mathrm{sep}} + \lambda_{\mathrm{bd}} \mathcal{E}_{\mathrm{bd}} + \lambda_{\mathrm{tr}} \mathcal{E}_{\mathrm{tr}}$\\[4pt]
$\downarrow$\\[4pt]
\textbf{최소화:} $u^* = \arg\min_{u \in \Sigma_m} \mathcal{E}(u)$ (PGD on $\Sigma_m$)\\[4pt]
$\downarrow$\\[4pt]
\textbf{진단:} $\mathbf{d} = (\mathsf{Bind}, \mathsf{Sep}, \mathsf{Inside}, \mathsf{Persist}) \in [0,1]^4$\\[4pt]
$\downarrow$\\[4pt]
\textbf{이산 복원:} Core, Interior, Boundary Band, Exterior
\vspace{0.5cm}}}
\caption{SCC 이론의 전체 파이프라인. 형식 우주의 여섯 구성 요소로부터 출발하여 공리적 제약, 에너지 최소화, 진단, 이산 복원으로 이어지는 계층적 구조. 각 단계는 Part~II--IV에서 상세히 전개된다.}
\label{fig:theory-pipeline}
\end{figure}


\subsection{진단 벡터: 형성의 질을 측정하다}

형성의 질을 평가하기 위해, SCC는 네 요소의 \textbf{proto-cohesion 진단 벡터}를 정의한다:
\begin{equation}
    \mathbf{d} = (\mathsf{Bind},\ \mathsf{Sep},\ \mathsf{Inside},\ \mathsf{Persist}) \in [0,1]^4
    \label{eq:diagnostic-vector-preview}
\end{equation}

각 요소는 $[0,1]$ 범위의 값으로, 형성이 각 구조적 요구를 얼마나 잘 만족하는지를 연속적으로 평가한다. 이 벡터의 모든 요소가 1에 가까운 형성은 완전한 객체에 가깝다. 일부만 높은 형성은 부분적으로 응집된 전-객체(pre-object)이다. 이에 대해서는 다음 장(Chapter~\ref{ch:diagnostic-vector})에서 상세히 다룬다.


\section{역사적 맥락과 관련 분야}
\label{sec:historical-context}

SCC 이론은 여러 지적 전통의 교차점에 위치한다.


\subsection{게슈탈트 심리학과 형태 이론}

게슈탈트 심리학의 핵심 통찰---``전체는 부분의 합 이상이다''---은 SCC의 직접적 선행자이다. Max Wertheimer, Wolfgang K\"ohler, Kurt Koffka가 1920년대에 정립한 게슈탈트 원리들---근접성(proximity), 유사성(similarity), 닫힘(closure), 연속성(continuity), 좋은 형태(Pr\"agnanz)---은 지각적 조직화가 개별 자극들의 단순 집합이 아니라 \emph{장(field)} 역학에 의해 결정된다는 관점을 제시했다. SCC의 응집 장은 이 직관의 수학적 실현이다.

그러나 SCC는 게슈탈트를 넘어선다. 게슈탈트 이론은 ``좋은 형태(good Gestalt)''의 원리를 정성적으로 기술했지만, 형태 형성의 역학적 과정에 대한 수학적 이론은 제시하지 못했다. SCC는 에너지 범함수의 최소화라는 변분 원리를 통해 형태 형성의 수학적 메커니즘을 명시적으로 제공한다. 구체적으로:
\begin{itemize}
    \item 게슈탈트의 ``근접성''은 SCC의 인접 핵 $\mathbf{N}_t$로 형식화된다.
    \item ``닫힘''은 닫힘 연산자 $\mathrm{Cl}_t$와 에너지 항 $\mathcal{E}_{\mathrm{cl}}$로 형식화된다.
    \item ``좋은 형태''는 에너지 최소화의 결과로서 자연스럽게 나타난다.
    \item ``전체는 부분의 합 이상''은 에너지의 비선형성과 자기-참조적 구조에서 구현된다.
\end{itemize}

% [Figure placeholder: historical timeline]
\begin{figure}[t]
\centering
\fbox{\parbox{0.85\textwidth}{\centering\vspace{3.5cm}
\textbf{[그림 1.5]} SCC 이론의 지적 계보.\\[4pt]
1920s: 게슈탈트 심리학 (Wertheimer, K\"ohler)\\
1945: Merleau-Ponty, \textit{지각의 현상학}---``전-객체적'' 개념\\
1958: Allen-Cahn 위상 장 방정식\\
1975: Kantorovich-Wasserstein 최적 전달\\
1993: 지속적 상동성 (Edelsbrunner et al.)\\
2000s: 비균형 최적 전달 (Chizat, Peyr\'e)\\
2025: SCC---다섯 전통의 통합
\vspace{0.5cm}}}
\caption{SCC 이론에 이르는 지적 계보. 다섯 가지 핵심 전통---게슈탈트 심리학, 현상학, 위상 장 이론, 최적 전달, 계산 위상수학---이 SCC에서 통합된다.}
\label{fig:intellectual-timeline}
\end{figure}


\subsection{위상 장 이론과 Allen-Cahn 방정식}

SCC의 에너지 구조에서 경계/형태학 항 $\mathcal{E}_{\mathrm{bd}}$는 Allen-Cahn 에너지와 구조적 유사성을 가진다:
\begin{equation}
    \mathcal{E}_{\mathrm{AC}} = \alpha \int |\nabla u|^2\, dx + \beta \int W(u)\, dx
    \label{eq:allen-cahn}
\end{equation}
여기서 $W(u) = u^2(1-u)^2$는 이중 우물 포텐셜이다.

그러나 SCC는 Allen-Cahn 이론이 아니며, 이로 환원될 수 없다. Allen-Cahn은 단일 에너지 항을 가지지만, SCC는 네 개의 독립적 항을 가진다. 특히 닫힘 항 $\mathcal{E}_{\mathrm{cl}}$과 분리 항 $\mathcal{E}_{\mathrm{sep}}$의 자기-참조적(self-referential) 구조---장이 ``무엇이 응집인가''를 정의하면서 동시에 그 기준에 의해 평가받는---는 표준 위상 장 이론에서는 나타나지 않는 고유한 특성이다.


\subsection{최적 전달 이론}

SCC의 시간적 전달 핵 $\mathbf{M}_{t \to s}$는 최적 전달(optimal transport) 이론과 연결된다. 전달 핵은 시각 $t$의 응집 측도 $u_t \cdot \mu_{X_t}$와 시각 $s$의 응집 측도 $u_s \cdot \mu_{X_s}$ 사이의 부분-확률적(sub-stochastic) 커플링이다. 이는 비균형 최적 전달(unbalanced optimal transport)---질량의 생성과 소멸이 허용되는---의 틀과 자연스럽게 연결된다.

특히 SCC의 전달 비용은 자기-참조적이다: 비용이 응집 장 자체에 의존한다. 이러한 자기-참조적 최적 전달 문제는 수학적으로 새로운 것이며, 존재성과 유일성에 대한 완전한 결과는 아직 열려 있다.


\subsection{현상학: Merleau-Ponty의 전-객체적 세계}

SCC의 ``전-객체적 응집''이라는 개념은 Maurice Merleau-Ponty의 현상학적 분석과 깊이 공명한다. Merleau-Ponty는 \textit{지각의 현상학}(1945)에서, 지각 경험이 개별 객체들의 모자이크가 아니라 배경(ground) 위에 떠오르는 형태(figure)의 \emph{역동적 장}이라고 주장했다.

그의 핵심 통찰은 ``전-객체적 세계(pre-objective world)''라는 개념이다: 의식적 객체 인식에 앞서, 지각 장은 이미 구조화되어 있되, 그 구조화가 아직 개별 객체의 분리를 완료하지 않은 상태가 있다. 이것은 정확히 SCC가 $[0,1]^4$ 진단 공간에서 $(1,1,1,1)$ 아닌 점들로 기술하고자 하는 상태이다.

SCC는 현상학적 기술에 \emph{수학적 정밀성}을 부여한다: 전-객체적 응집의 정도를 네 차원에서 정량적으로 측정하고, 그 역학을 에너지 최소화로 기술하며, 객체의 창발을 위상 전이(phase transition)로 특성화한다.


\subsection{인지과학: 전-객체적 인식}

인지과학에서 ``전주의적(pre-attentive)'' 처리의 존재는 잘 알려져 있다. 시각 정보는 의식적 객체 인식에 앞서 전-의식적 수준에서 조직화된다. SCC의 ``전-객체적 응집(pre-objective cohesion)''은 이 전주의적 조직화의 수학적 모델로 해석될 수 있다.

SCC는 또한 주의(attention)와 의식의 통합 이론에도 시사점을 가진다. 형성의 네 진단 요소---결합, 분리, 내부성, 지속성---는 의식적 객체 인식의 필요조건들로 해석될 수 있으며, 진단 벡터의 연속적 특성은 ``의식의 문턱'' 같은 이분법적 모델에 대한 대안을 제시한다.


\section{이 책의 구조}
\label{sec:book-roadmap}

이 책은 SCC 이론을 체계적으로 전개한다.

\textbf{Part~I (기초)}에서는 이론의 동기(이 장)와 핵심 진단 개념(Chapter~\ref{ch:diagnostic-vector})을 소개한다. 수학적 엄밀성보다는 직관적 이해에 중점을 둔다.

\textbf{Part~II (수학적 구조)}에서는 에너지 범함수, 연산자 공리, 변분 구조를 수학적으로 엄밀하게 전개한다. 위상 전이 정리(T8), 기울기 흐름 수렴(T14), $\Gamma$-수렴(T11) 등 핵심 정리들을 증명한다.

\textbf{Part~III (시간적 구조와 다중 형성)}에서는 시간적 전달과 지속성의 이론, 그리고 여러 형성이 공존하는 다중 형성 체계를 다룬다.

\textbf{Part~IV (구현과 실험)}에서는 Python 구현체 \texttt{scc} 패키지의 구조와 실험 결과를 제시한다.

\textbf{Part~V (응용과 전망)}에서는 인지과학적 해석, 열린 문제, 미래 방향을 논의한다.

\bigskip

다음 장에서는 SCC 이론의 핵심 진단 도구인 proto-cohesion 진단 벡터 $\mathbf{d} = (\mathsf{Bind}, \mathsf{Sep}, \mathsf{Inside}, \mathsf{Persist})$의 각 요소를 직관과 구체적 예제를 통해 소개한다.


\section{토론 질문과 연습 문제}
\label{sec:ch1-exercises}

\subsection*{토론 질문}

\begin{enumerate}
    \item 일상에서 경험하는 ``경계가 모호한 형성''의 예를 세 가지 더 들어보라. 각각에서 소속과 비소속 사이의 전이가 어떻게 일어나는지 기술하라.

    \item ``분할의 문제''에서 논리적 순환을 구체적으로 설명하라: 영상 분할 알고리즘이 ``물건''을 찾기 위해 이미 ``물건''에 대한 어떤 지식을 전제해야 하는가?

    \item 퍼지 집합(fuzzy set) 이론과 SCC의 응집 장은 둘 다 $[0,1]$-값의 소속 함수를 사용한다. 그러나 \S\ref{sec:desiderata}의 다섯 가지 필요조건 중 퍼지 집합이 충족하지 못하는 것은 무엇인가?

    \item SCC의 ``존재론적 우선성 원리''에 따르면, 이산적 객체는 항상 파생적이다. 이 원리가 일상적 직관(``물건은 명확하게 존재한다'')과 어떻게 양립하는지 논의하라.

    \item 게슈탈트 심리학의 ``좋은 형태(Pr\"agnanz)'' 원리와 SCC의 에너지 최소화 원리를 비교하라. 어떤 면에서 SCC가 게슈탈트를 ``넘어서는가''?

    \item Merleau-Ponty의 ``전-객체적 세계'' 개념이 SCC에서 어떻게 수학적으로 구현되는지 설명하라. $[0,1]^4$ 진단 공간의 어떤 영역이 ``전-객체적'' 상태에 해당하는가?
\end{enumerate}


\subsection*{연습 문제}

\begin{enumerate}
    \item \textbf{(프로그래밍)} Python \texttt{scc} 패키지를 이용하여 $10 \times 10$ 격자 위에서 형성을 찾아보라:
    \begin{verbatim}
    from scc import GraphState, ParameterRegistry, find_formation
    g = GraphState.grid_2d(10, 10)
    p = ParameterRegistry()
    r = find_formation(g, p)
    print(r.diagnostics)
    \end{verbatim}
    결과의 네 진단 값을 해석하라.

    \item \textbf{(비교 분석)} $10 \times 10$ 격자 위에서 다음 세 가지 장을 구성하고, 각각의 진단 벡터를 비교하라: (a) Gaussian bump ($\sigma = 1.5$), (b) 균일 장 $u \equiv 0.3$, (c) 무작위 장. 표~\ref{tab:framework-comparison}의 결과와 연결지어 논의하라.

    \item \textbf{(개념적)} Allen-Cahn 에너지 $\mathcal{E}_{\mathrm{AC}} = \alpha \int |\nabla u|^2 + \beta \int W(u)$와 SCC 에너지 $\mathcal{E} = \lambda_{\mathrm{cl}} \mathcal{E}_{\mathrm{cl}} + \lambda_{\mathrm{sep}} \mathcal{E}_{\mathrm{sep}} + \lambda_{\mathrm{bd}} \mathcal{E}_{\mathrm{bd}} + \lambda_{\mathrm{tr}} \mathcal{E}_{\mathrm{tr}}$를 비교하라. SCC에서 Allen-Cahn에 해당하는 항은 무엇인가? SCC에만 존재하는 항은 무엇이며, 그 역할은 무엇인가?
\end{enumerate}

\begin{center}
\fbox{\parbox{0.85\textwidth}{\centering
\textbf{교수자를 위한 노트}\\[6pt]
이 장은 수학적 엄밀성보다 직관적 동기 부여에 중점을 둔다.\\
토론 질문 1--4는 철학적/인지과학적 배경의 학생에게 적합하고,\\
질문 5--6과 연습 문제 1--3은 수학/공학 배경의 학생에게 적합하다.\\
연습 문제 1은 \texttt{scc} 패키지 설치가 필요하다 (부록 참조).
}}
\end{center}
